\section*{Chapitre \ref{ch:b2:isotope-geochemistry} --- Géochimie isotopique}

\begin{solution}{exo:b2:isotope-geochemistry:1}
Tristan~: $(0.51249/0.512638 - 1)\times10^4 = -2.9$. Mangaia~: $(0.51285/0.512638 -
1)\times10^4 = +4.1$. MORB~: $+8.5$. Les sources de Mangaia et du MORB, à $\varepsilon_{\mathrm{Nd}}$ positif,
ont eu en moyenne un Sm/Nd plus élevé que les \omterm{def:b2:planet-formation-meteorites:chondrule}{chondrites}~; celle de Tristan en a eu un
plus bas.
\end{solution}

\begin{solution}{exo:b2:isotope-geochemistry:2}
$\delta^{18}\mathrm{O}_{\mathrm{VPDB}} = (\num{2.0672e-3}/\num{2.0052e-3} - 1)\times1000 =
+30.9$ pour mille sur l'échelle \omterm{def:b2:isotope-geochemistry:standards}{VSMOW}. La calcite a $^{18}\mathrm{O}/^{16}\mathrm{O} =
\num{2.0672e-3}\times(1 - 0.0020) = \num{2.0631e-3}$, d'où $\delta^{18}\mathrm{O} =
(\num{2.0631e-3}/\num{2.0052e-3} - 1)\times1000 = +28.9$ pour mille (\omterm{def:b2:isotope-geochemistry:standards}{VSMOW}).
\end{solution}

\begin{solution}{exo:b2:isotope-geochemistry:3}
À \qty{1.0}{Ga}~: $e^{\lambda t} - 1 = 0.006545$, et $0.512638 - 0.1966\times0.006545 =
0.511351$. À \qty{4.567}{Ga}~: $e^{\lambda t}-1 = 0.030243$, et $0.512638 -
0.1966\times0.030243 = 0.506692$.
\end{solution}

\begin{solution}{exo:b2:isotope-geochemistry:4}
$\delta\mathrm{D} = 8\times(-18.0) + 10 = -134$ pour mille. Pour le lac, $d = -42 -
8\times(-4.0) = -10$ pour mille, en dessous des 10 de la droite~: l'eau se situe sous la \omterm{prop:b2:isotope-geochemistry:mwl}{droite
des eaux météoriques}, comme le font les eaux de surface évaporées (les isotopes légers sont partis en premier).
\end{solution}

\begin{solution}{exo:b2:isotope-geochemistry:5}
$Q = 10^4\times\num{6.524e-12}\times0.1966/0.512638 = \num{2.502e-8}$ par an, soit
\qty{25.0}{\per\Ga}. Le granite a $\mathrm{Sm}/\mathrm{Nd} = 6.0/30 = 0.200$,
contre $0.145/0.457 = 0.317$ pour le CI~: $f = 0.200/0.317 - 1 = -0.37$. Sur
\qty{0.5}{Ga}~: $\Delta\varepsilon = 25.0\times(-0.37)\times0.5 = -4.6$ unités.
\end{solution}

\begin{solution}{exo:b2:isotope-geochemistry:6}
$\varepsilon \approx 3.5 + 25.0\times(-0.45)\times2.7 = 3.5 - 30.4 = -26.9$, contre
$-27.3$~: l'approximation est trop élevée de 0.4 unité (elle n'est pas assez négative). L'essentiel de
l'écart, environ 0.3 unité, provient de $e^{\lambda t} > 1$~: le parent, donc le
déficit de croissance relatif au CHUR, était plus abondant dans le passé qu'aujourd'hui
(voir l'\cref{exo:b2:isotope-geochemistry:10})~; le reste vient de l'arrondi de $f$ et de la valeur initiale.
\end{solution}

\begin{solution}{exo:b2:isotope-geochemistry:7}
Avec une fraction massique $x$ de granite, $\mathrm{Sr} = 320x + 130(1-x)$ et
$R = (320x\times0.7200 + 130(1-x)\times0.7030)/\mathrm{Sr}$. Pour $x = 0.05$~: Sr
\qty{139.5}{ppm}, $R = 0.70495$~; $x = 0.10$~: 149.0, 0.70665~; $x = 0.20$~: 168.0,
0.70948. D'après la démonstration de la proposition de mélange, $R\,\mathrm{Sr}$ est linéaire en Sr~:
$R = 0.731632 - 3.7221/\mathrm{Sr}$, une droite de pente $-3.722$ (ppm) et
d'ordonnée à l'origine 0.7316 en fonction de $1/\mathrm{Sr}$. Vérification pour $x = 0.05$~: $0.731632 - 3.7221/139.5
= 0.70495$.
\end{solution}

\begin{solution}{exo:b2:isotope-geochemistry:8}
À \qty{20}{\degreeCelsius}~: $\delta^{18}\mathrm{O} = 1000/1.0098 - 1000 = -9.70$ pour
mille, $\delta\mathrm{D} = 1000/1.084 - 1000 = -77.5$ pour mille, pente $77.5/9.70 = 7.99$.
À \qty{0}{\degreeCelsius}~: $-11.56$ et $-99.9$ pour mille, pente 8.64. Les deux sont proches de
la pente 8 de la \omterm{prop:b2:isotope-geochemistry:mwl}{droite des eaux météoriques}~: le seul fractionnement d'équilibre de l'eau
entre liquide et vapeur fixe le rapport des décalages du deutérium et de l'oxygène 18
voisin de 8.
\end{solution}

\begin{solution}{exo:b2:isotope-geochemistry:9}
La croûte contient $0.0051\times27/1.25 = 0.110$, soit 11 pour cent du néodyme. Bilan de
masse~: $0.110\times(-15) + 0.890\,\varepsilon_m = 0$, donc $\varepsilon_m = 1.65/0.890 =
+1.9$. Avec un réservoir appauvri détenant une part $x$ à $+8.5$ et le reste à 0~:
$0.110\times(-15) + 8.5x = 0$, $x = 1.65/8.5 = 0.19$~: environ un cinquième du
néodyme.
\end{solution}

\begin{solution}{exo:b2:isotope-geochemistry:10}
Le taux de croissance vaut $\lambda p\,e^{\lambda t}$, avec $e^{\lambda t} \approx 1 + \lambda t$~;
moyenné sur l'intervalle allant de la séparation $T$ au présent, le facteur vaut
$1 + \lambda T/2$, donc les termes négligés sont de taille relative $\lambda T/2$. Sur
l'âge de la Terre, $\num{6.524e-12}\times\num{4.567e9}/2 = 0.015$, soit 1.5 pour cent. Pour la
croûte de l'\cref{exo:b2:isotope-geochemistry:6}, $\lambda T/2 = 0.0088$, et l'erreur vaut $0.0088\times30.4 = 0.27$
unité. La lente variation du rapport CHUR, par laquelle $\varepsilon$ est divisé, ajoute
un autre terme, plus petit, d'ordre $(p/R)\lambda T \times \varepsilon_0$, soit ici 0.02 unité.
\end{solution}

\begin{solution}{exo:b2:isotope-geochemistry:11}
$K = 13.3/5.0 = 2.66$. Avec $s = 0.5$~: $s' = 2.66\times0.5/(1 + 1.66\times0.5) =
1.33/1.83 = 0.73$. L'écart $s' - s$ est maximal là où
$\mathrm{d}s'/\mathrm{d}s = K/(1+(K-1)s)^2 = 1$, soit en $s = (\sqrt K - 1)/(K-1) =
0.631/1.66 = 0.38$, où $s' = 0.62$~: écart de 0.24.
\end{solution}

\begin{solution}{exo:b2:isotope-geochemistry:12}
Production~: $2.3$ à $2.6\times\qty{7.2e9}{km^3}/\qty{4.5e9}{an} = 3.7$ à
\qty{4.2}{km^3} par an. Destruction~: $1.3$ à $1.6\times7.2/4.5 = 2.1$ à
\qty{2.6}{km^3} par an. Les taux actuels diffèrent~: les apports (2.7) sont inférieurs à la
production moyenne du modèle, le retrait (3.2) supérieur à sa destruction moyenne. Dans le modèle, la croûte a grandi
tôt, quand la production était la plus rapide, et le présent est une époque d'équilibre ou de légère
perte, comme le dit le bilan des zones de subduction.
\end{solution}

\begin{solution}{pb:b2:isotope-geochemistry:1}
\textbf{1.} Manteau~: $(0.51307/0.512638 - 1)\times10^4 = +8.4$. Sédiment~:
$(0.51141/0.512638 - 1)\times10^4 = -24.0$.

\textbf{2.} Kora $+2.8$, Vel $-4.3$, Asha $+4.7$.

\textbf{3.} Avec $^{87}\mathrm{Sr}/^{86}\mathrm{Sr}$ en abscisse et $\varepsilon_{\mathrm{Nd}}$ en ordonnée~:
manteau (0.70282, $+8.4$), Kora (0.70410, $+2.8$), Vel (0.70618, $-4.3$), sédiment
(0.72000, $-24.0$), Asha (0.70290, $+4.7$). Seule Vel, avec un $\varepsilon_{\mathrm{Nd}}$ négatif,
se situe du côté enrichi du CHUR.

\textbf{4.} Sous la bande mantellique~: au rapport du strontium du \omterm{def:b2:isotope-geochemistry:reservoirs}{manteau appauvri}, la bande
aurait un $\varepsilon_{\mathrm{Nd}}$ voisin de $+8$~; Asha a $+4.7$. Il ressemble au
groupe de Sainte-Hélène de type HIMU (Mangaia, Tubuai, Sainte-Hélène).

\textbf{5.} La fusion modifie les rapports parent--fils, mais pas, de façon mesurable, les
rapports isotopiques du strontium et du néodyme (proposition sur la fusion)~; les laves sont
jeunes, de sorte que leurs rapports n'ont pas augmenté depuis l'éruption (question 6).

\textbf{6.} $\Delta R = 0.12\,(e^{\lambda\cdot10^6} - 1) = 0.12\times\num{6.52e-6} =
\num{7.8e-7}$, soit $\num{7.8e-7}/0.512638\times10^4 = 0.015$ unité epsilon~: bien en dessous des
différences entre les îles, et en dessous de la précision d'une mesure.

\textbf{7.} $\mathrm{Nd} = 40f + 0.75(1-f)$, et
$R = \bigl(40f\times0.51141 + 0.75(1-f)\times0.51307\bigr)/\mathrm{Nd}$.

\textbf{8.} $f = 0.5$ pour cent~: Nd \qty{0.9463}{ppm}, $\varepsilon = +1.6$~; 1 pour cent~:
1.1425, $-2.9$~; 2 pour cent~: 1.535, $-8.4$~; 5 pour cent~: 2.7125, $-15.4$.

\textbf{9.} Sr 10.95, 11.90, 13.80 et \qty{19.50}{ppm}~; $^{87}\mathrm{Sr}/^{86}\mathrm{Sr} =
0.70439$, 0.70571, 0.70780 et 0.71163.

\textbf{10.} $R\,\mathrm{Sr} = 10\times0.70282 + f(200\times0.72000 - 10\times0.70282)$ et
$\mathrm{Sr} = 10 + 190f$~; en éliminant $f$, $R\,\mathrm{Sr} = 7.0282 +
(\mathrm{Sr}-10)\times136.9718/190$, donc $R = 0.720904 - 0.18084/\mathrm{Sr}$~: une droite
de pente $-0.1808$ (ppm) en fonction de $1/\mathrm{Sr}$. Vérification pour $f = 1$ pour cent~:
$0.720904 - 0.18084/11.90 = 0.70571$.

\textbf{11.} $K = (10/0.75)/(200/40) = 13.33/5.0 = 2.67$. Le manteau est relativement
plus riche en strontium, le sédiment en néodyme, si bien que le rapport du néodyme se déplace
vers celui du sédiment plus vite~: un peu de sédiment déplace $\varepsilon_{\mathrm{Nd}}$ bien avant que
$^{87}\mathrm{Sr}/^{86}\mathrm{Sr}$ ne bouge beaucoup.

\textbf{12.} 53 fois plus riche en Nd ($40/0.75$), 20 fois en Sr. À $f = 1$ pour cent le
sédiment fournit $0.40/1.1425 = 35$ pour cent du néodyme, ce qui déplace
$\varepsilon_{\mathrm{Nd}}$ de $0.35\times(8.4 + 24.0) = 11.3$ unités, mais seulement $2.0/11.9 =
17$ pour cent du strontium.

\textbf{13.} $0.01\times40/1.1425 = 0.350$~: 35 pour cent.

\textbf{14.} À partir de $R\,(40f + 0.75(1-f)) = 40fR_S + 0.75(1-f)R_D$~:
$f = 0.75(R - R_D)/\bigl(40(R_S - R) + 0.75(R - R_D)\bigr)$, et de même avec 200
et 10 pour le strontium.

\textbf{15.} $f = 0.75\times(-0.00029)/\bigl(40\times(-0.00137) + 0.75\times(-0.00029)\bigr)
= \num{-2.175e-4}/\num{-5.502e-2} = 0.0040$~: 0.40 pour cent.

\textbf{16.} $f = 10\times0.00128/\bigl(200\times0.01590 + 10\times0.00128\bigr) =
0.0128/3.1928 = 0.0040$~: 0.40 pour cent~; les deux concordent.

\textbf{17.} Néodyme~: $0.75\times(-0.00065)/\bigl(40\times(-0.00101) +
0.75\times(-0.00065)\bigr) = \num{-4.875e-4}/\num{-4.089e-2} = 0.0119$~; strontium~:
$10\times0.00336/\bigl(200\times0.01382 + 10\times0.00336\bigr) = 0.0336/2.7976 = 0.0120$.
Les deux donnent 1.2 pour cent.

\textbf{18.} $f = 0.75\times(-0.00019)/\bigl(40\times(-0.00147) + 0.75\times(-0.00019)\bigr)
= 0.0024$, ce qui prédit Sr \qty{10.46}{ppm} et $^{87}\mathrm{Sr}/^{86}\mathrm{Sr} =
0.70361$, contre 0.70290 mesuré. Asha ne se situe pas sur la courbe de mélange manteau--sédiment~:
sa source contient un autre constituant.

\textbf{19.} Vieille croûte océanique recyclée, de type HIMU~: ses isotopes du néodyme demandent un Sm/Nd
un peu inférieur à celui du \omterm{def:b2:isotope-geochemistry:reservoirs}{manteau appauvri}, ses isotopes du strontium un Rb/Sr faible, et son
plomb (non mesuré ici) un U/Pb élevé, le tout sur un à deux milliards d'années.

\textbf{20.} $e^{\lambda_{147}\cdot2\times10^9} - 1 = 0.013133$. Sédiment~: $0.51141 -
0.110\times0.013133 = 0.509965$~; CHUR~: $0.512638 - 0.1966\times0.013133 = 0.510056$~;
$\varepsilon_{\mathrm{Nd}}(2\,\mathrm{Ga}) = (0.509965/0.510056 - 1)\times10^4 = -1.8$.

\textbf{21.} $f = 0.110/0.1966 - 1 = -0.440$~; depuis la subduction, $\varepsilon$ a varié de
$25.0\times(-0.440)\times2.0 = -22.0$, donc $\varepsilon(2\,\mathrm{Ga}) \approx -24.0 + 22.0
= -2.0$, à 0.2 unité près de la valeur exacte.

\textbf{22.} $e^{\lambda_{87}\cdot2\times10^9} - 1 = 0.028338$~;
$0.72000 - 0.50\times0.028338 = 0.70583$.

\textbf{23.} Au moment de la subduction, le sédiment avait $\varepsilon_{\mathrm{Nd}} = -1.8$, proche de
la valeur $-4.7$ d'une croûte âgée de 0.7 milliard d'années seulement à cette époque, et nettement au-dessus des
valeurs fortement négatives de la vieille croûte~: cela s'accorde avec l'érosion d'une croûte jeune, peut-être
mêlée d'un peu de matériau issu du manteau. Son $^{87}\mathrm{Sr}/^{86}\mathrm{Sr}$ de
0.7058 à cette époque est aussi celui d'une roche crustale jeune.

\textbf{24.} La source de Vel, l'île la plus enrichie, contient
\textbf{\qty[text-series-to-math=true]{1.2}{\percent}} de sédiment recyclé.
\end{solution}
